\chapter{Cross-References and Citations}\label{chap:references}

This chapter covers the second half of the fourth block:
referring to sections, equations, figures and tables by their number,
and citing sources with a bibliography that builds itself.

\section{Labels}\label{sec:labels}

Anything that \hologo{LaTeX} numbers can get a label, a name to refer to it:

\slidecode{04-labels.tex}

\begin{itemize}
    \item A section gets its label on the same line as \code|\section|.
        With a paragraph break in between, \code|\pageref| can point at the wrong page.
    \item An equation gets its label inside the environment.
        In \texttt{align}, each numbered line can have its own.
    \item A table or a figure gets its label after \code|\caption| (\autoref{sec:floats}).
\end{itemize}

The name is up to you, but it must be unique in the document;
a name used twice gives the warning \texttt{Label `...' multiply defined},
and references to it are unreliable.
A prefix for the kind of thing makes a long list of labels readable,
and the editor's autocompletion groups them:

\begin{description}
    \item[\texttt{sec:}] sections (and \texttt{chap:} for chapters)
    \item[\texttt{eq:}] equations
    \item[\texttt{fig:}] figures
    \item[\texttt{tab:}] tables
\end{description}

\hologo{LaTeX} itself ignores the prefixes; \texttt{fig:results} is just a name.

\section{References}\label{sec:refs}

\code|\ref| prints the number of a label, \code|\eqref| the number of an equation in brackets:

\slidecode{04-ref.tex}
\sliderender{04-ref-render.pdf}

\begin{itemize}
    \item \code|\ref| only gives the number; the word in front, \enquote{Figure}, is typed by you.
        \autoref{sec:hyperref} gets rid of that too.
    \item \code|\eqref| comes from \texttt{amsmath}.
        It gives \enquote{(1)}, the usual way to refer to an equation.
    \item The tilde is the non-breaking space from \autoref{sec:spaces}.
        Without it, \enquote{Figure} can end one line and \enquote{1} start the next.
    \item \code|\pageref| gives the page instead of the number.
        Write it with a tilde too: \code|page~\pageref{fig:cooling}|.
\end{itemize}

No number is typed by hand.
Move the figure, add an equation, and every number and every reference follows,
like the table of contents.
Like the table of contents, references need two compilations:
the first one collects the numbers, the second one prints them (\autoref{sec:compiling}).
A reference to a label that does not exist prints \texttt{??} (\autoref{sec:undefined-refs}).

\section{Automatic names: \texttt{hyperref}}\label{sec:hyperref}

The package \texttt{hyperref} adds \code|\autoref|, which writes the name in front of the number itself:

\slidecodetitled{04-autoref-preamble.tex}{Preamble}
\slidecodetitled{04-autoref.tex}{Document}
\sliderender{04-autoref-render.pdf}

\begin{itemize}
    \item The name comes from what was labelled:
        a \texttt{figure} gives \enquote{Figure}, a \texttt{table} \enquote{Table}, a \code|\section| \enquote{section}.
        With \texttt{babel}, the names follow the language of the document,
        e.g.\ \enquote{Abbildung} in German (\autoref{sec:babel}).
    \item \code|\autoref| writes the non-breaking space itself.
    \item Equations come out as \enquote{Equation 1}, without brackets.
        For equations, \code|Equation~\eqref{...}| stays the better choice.
\end{itemize}

\texttt{hyperref} also turns every reference, every citation and the table of contents into a link,
fills in the bookmarks of the PDF,
and adds \code|\href{https://ctan.org}{CTAN}| (a link with its own text)
and \code|\url{https://ctan.org}| (an address, printed as it is) for links to websites.

\subsection{Load it last}

\texttt{hyperref} changes many other packages, so it goes after them, at the end of the preamble.
This is the one package where the order matters
(the few packages that must come after it say so in their manual; \texttt{cleveref} is one of them).

\subsection{Coloured boxes around links}

By default, links get coloured boxes in a PDF viewer.
They are only shown on screen, never printed, but most people find them ugly.
\code|\hypersetup{hidelinks}|, after \code|\usepackage{hyperref}|, removes them.
The links still work.

\subsection{Where the reference commands come from}

\autoref{tab:ref-commands} lists the reference commands and the package each one needs.
\code|\nameref| prints the title of a section instead of its number:

\owncode{ch07-nameref.tex}
\ownrender{ch07-nameref.pdf}

\begin{table}[htbp]
    \centering
    \caption{Reference commands and where they come from.}
    \label{tab:ref-commands}
    \begin{tabular}{@{}lll@{}}
        \toprule
        \textbf{Command} & \textbf{Prints} & \textbf{Package} \\
        \midrule
        \code|\ref|       & the number                       & \hologo{LaTeX} itself \\
        \code|\pageref|   & the page                         & \hologo{LaTeX} itself \\
        \code|\eqref|     & the number in brackets            & \texttt{amsmath} \\
        \code|\autoref|   & name and number                   & \texttt{hyperref} \\
        \code|\nameref|   & the title (of a section, or the caption) & \texttt{hyperref} \\
        \code|\cref|      & name and number, also several at once & \texttt{cleveref} \\
        \bottomrule
    \end{tabular}
\end{table}

\section{\texttt{cleveref}}\label{sec:cleveref}

The package \texttt{cleveref} does what \code|\autoref| does, with more control.
It is loaded after \texttt{hyperref}:

\owncodetitled{ch07-cleveref-preamble.tex}{Preamble}
\owncodetitled{ch07-cleveref.tex}{Document}
\ownrender{ch07-cleveref-render.pdf}

\begin{itemize}
    \item \code|\cref| writes the name and the number, like \code|\autoref|.
        Equations get their brackets.
    \item Several labels in one \code|\cref| are combined and sorted: \enquote{Figures 1 to 3}.
    \item Without the option \texttt{noabbrev}, names are abbreviated (\enquote{fig.}, \enquote{eq.}).
        Without \texttt{capitalise}, they are lower case in the middle of a sentence,
        and \code|\Cref| capitalises them at the start of one.
    \item \code|\crefname{figure}{Fig.}{Figs.}| in the preamble changes a name.
\end{itemize}

For a lab report, \code|\autoref| is enough.
\texttt{cleveref} is worth it in a thesis, where references to several figures or to ranges of equations are common.

\section{The bibliography file}\label{sec:bibfile}

Every source is written down once, in a file of its own next to the document, usually \texttt{references.bib}.
Each source is an entry:

\slidecodetitled[language=bib]{04-bib-entry.bib}{\texttt{references.bib}}

\begin{description}
    \item[Entry type] after the \code|@|: what kind of source it is (\autoref{tab:entry-types}).
    \item[Key] your choice, the name you cite it by.
        Author plus year is the usual pattern, e.g.\ \texttt{halliday2018}.
    \item[Fields] \code|name = {value}|, separated by commas.
        The order does not matter; a comma after the last one is allowed.
\end{description}

The format is called Bib\hologo{TeX}, after the program that introduced it in the 1980s.
It is still the format that every reference manager and publisher website exports.

\subsection{Entry types}

\autoref{tab:entry-types} lists the common entry types and the fields \texttt{biblatex} needs for each
\parencite{biblatex}.
More fields are always allowed: \texttt{doi}, \texttt{isbn}, \texttt{url}, \texttt{pages}, \texttt{edition}, \texttt{publisher}.
A missing field is not an error, but the reference list will lack that information.

\begin{table}[htbp]
    \centering
    \caption{Common entry types and the fields they need.}
    \label{tab:entry-types}
    \begin{tabularx}{\linewidth}{@{}l>{\raggedright\arraybackslash}X>{\raggedright\arraybackslash}X@{}}
        \toprule
        \textbf{Type} & \textbf{For} & \textbf{Fields} \\
        \midrule
        \texttt{@article}       & a journal article     & author, title, journaltitle, year; often volume, number, pages, doi \\
        \texttt{@book}          & a book                & author, title, year; often publisher, edition \\
        \texttt{@incollection}  & a chapter by its own authors in an edited book & author, title, booktitle, year; often editor, publisher, pages \\
        \texttt{@inproceedings} & a conference paper    & author, title, booktitle, year \\
        \texttt{@thesis}        & a thesis              & author, title, type (e.g.\ {Bachelor's thesis}), institution, year \\
        \texttt{@report}        & a technical report    & author, title, type, institution, year \\
        \texttt{@manual}        & a manual, a datasheet, a standard & author or organisation, title, year \\
        \texttt{@online}        & a web page            & author or organisation, title, url, year; urldate (when you read it) \\
        \texttt{@misc}          & anything else         & author, title, year \\
        \bottomrule
    \end{tabularx}
\end{table}

Exports often use the older names of Bib\hologo{TeX}:
\texttt{journal} instead of \texttt{journaltitle},
\texttt{@phdthesis}, \texttt{@mastersthesis} and \texttt{@techreport}.
\texttt{biblatex} understands all of them.

\subsection{Author names}

\autoref{tab:author-names} shows how to write the \texttt{author} field.
Bib\hologo{TeX} splits every name into parts (given name, family name and a few more),
and the citation style decides how to put them back together.

\begin{table}[htbp]
    \centering
    \caption{Names in the \texttt{author} field.}
    \label{tab:author-names}
    \begin{tabular}{@{}ll@{}}
        \toprule
        \textbf{In the \texttt{author} field} & \textbf{Meaning} \\
        \midrule
        \code|Halliday, David|                       & family, given \\
        \code|David Halliday|                        & also works \\
        \code|Halliday, David and Resnick, Robert|   & several authors: \code|and| \\
        \code|Halliday, David and others|            & \enquote{et al.} \\
        \code|van Beethoven, Ludwig|                 & \enquote{van} stays with the family name \\
        \code|{{World Health Organization}}|         & organisation: one name \\
        \bottomrule
    \end{tabular}
\end{table}

\begin{itemize}
    \item With a comma, everything before it is the family name.
        Without a comma, the last word is.
        Both give the same result; exports use the comma form.
    \item Names are separated by \code|and|, never by commas.
        \texttt{Halliday, David, Resnick, Robert} is read as one strange name.
        This is the mistake you are most likely to make when typing an entry by hand.
    \item Prefixes such as \enquote{van} or \enquote{von} are recognised because they are lower case.
    \item An organisation needs a second pair of braces, otherwise its last word is taken as a family name.
\end{itemize}

The last rule is easy to miss, because it only shows with some styles.
Cited with \code|\textcite| and \texttt{style=authoryear} (\autoref{sec:cite-styles}),
without the inner braces:

\slidecodetitled[language=bib]{04-org-wrong.bib}{\texttt{references.bib}}
\sliderender{04-org-wrong.pdf}

\enquote{World Health} is taken as the given name, \enquote{Organization} as the family name.
With double braces, the outer pair belongs to the field and the inner pair makes the whole name one unit:

\slidecodetitled[language=bib]{04-org-right.bib}{\texttt{references.bib}}
\sliderender{04-org-right.pdf}

\subsection{Capitals in titles}

Some styles change titles to sentence case (\texttt{apa}, for example),
so \enquote{Cooling of DNA Samples} would become \enquote{Cooling of dna samples}.
Braces protect what must keep its capitals: \code|title = {Cooling of {DNA} Samples}|.
The standard \texttt{biblatex} styles leave titles as they are,
so with them the braces are not needed, but they do no harm.

\subsection{Finding entries}

Nobody types these entries by hand.
Copy them, but check them:

\begin{itemize}
    \item \textbf{Google Scholar}: the quotation mark icon under a result, then \textit{BibTeX} at the bottom.
    \item \textbf{arXiv}: \textit{Export BibTeX Citation} in the right-hand column of the abstract page.
        If the paper has since been published, cite the published version instead.
    \item \textbf{Publisher or library catalogue}: \textit{Export citation}, with a Bib\hologo{TeX} option.
    \item \textbf{From a DOI}: \myhref{https://www.doi2bib.org}{doi2bib.org}.
        The \gls{doi} is the identifier starting with \texttt{10.} printed on almost every paper.
    \item \textbf{A reference manager}: Zotero (with the Better BibTeX plug-in) or JabRef
        keep the \texttt{.bib} file up to date for you,
        and add entries from a DOI, an ISBN or an arXiv ID.
        Worth it as soon as you write a thesis.
\end{itemize}

Exported entries are often incomplete or oddly capitalised,
organisations come without the double braces,
and the reference list only shows what is in the entry.

\section{\texttt{biblatex} and \texttt{biber}}\label{sec:biber}

The package \texttt{biblatex} reads the \texttt{.bib} file and formats the citations and the reference list:

\slidecodetitled{04-biblatex-preamble.tex}{Preamble}
\slidecodetitled{04-cite.tex}{Document}
\sliderender{04-cite.pdf}

\begin{itemize}
    \item \code|\addbibresource| takes the file name \emph{with} \texttt{.bib}.
    \item \code|\cite{key}| cites an entry (more commands in \autoref{sec:cite-commands}).
    \item \code|\printbibliography| sets the reference list, with its own heading.
        Only cited entries appear in it.
        \code|\nocite{*}| adds all entries of the file, \code|\nocite{key}| a single one without citing it.
    \item \code|\printbibliography[heading=bibintoc]| also adds it to the table of contents.
\end{itemize}

\subsection{Two parts that work together}

The names sound like two competing tools, but they are two halves of one setup:

\begin{description}
    \item[\texttt{biblatex}] a \hologo{LaTeX} package: formats the citations and the list.
    \item[\texttt{biber}] a program: reads the \texttt{.bib} file and sorts the entries.
\end{description}

\texttt{biber} runs between two \hologo{LaTeX} runs:
\hologo{LaTeX} writes down which keys were cited,
\texttt{biber} looks them up in the \texttt{.bib} file and sorts them,
and the next \hologo{LaTeX} run sets the result.
Overleaf and \texttt{latexmk} notice when this is needed and run \texttt{biber} for you.
By hand, the sequence is \texttt{pdflatex}, \texttt{biber}, \texttt{pdflatex}, \texttt{pdflatex} (\autoref{sec:compiling}).
\code|backend=biber| is the default; it is written out so it is clear which program runs.

\texttt{biblatex} with \texttt{babel} warns that it would like the package \texttt{csquotes} too;
load it, and the warning disappears (\autoref{sec:more-packages}).

\subsection{Older setups}

Older tutorials and templates use a different setup:
the package \texttt{natbib}, or \code|\bibliographystyle{plain}| and \code|\bibliography{references}|,
together with the program \texttt{bibtex} instead of \texttt{biber}.
It still works, but it has fewer styles and options, and trouble with umlauts and accents.
\texttt{biblatex} can also use the old program (\code|backend=bibtex|), with the same limits.
For a new document, use \texttt{biblatex} with \texttt{biber};
if your department's template uses the older setup, stay with it.

\section{Citation commands}\label{sec:cite-commands}

\autoref{tab:cite-commands} lists the citation commands of \texttt{biblatex}.

\begin{table}[htbp]
    \centering
    \caption{Citation commands of \texttt{biblatex}.}
    \label{tab:cite-commands}
    \begin{tabular}{@{}ll@{}}
        \toprule
        \textbf{Command} & \textbf{Use} \\
        \midrule
        \code|\textcite{key}|          & as part of the sentence: \enquote{Halliday et al.\ [1] show} \\
        \code|\parencite{key}|         & in brackets, as an aside \\
        \code|\parencite[p.~12]{key}|  & with a page (or chapter, figure, \ldots) \\
        \code|\parencite{key1,key2}|   & several sources at once \\
        \code|\cite{key}|              & depends on the style \\
        \code|\footcite{key}|          & as a footnote \\
        \code|\citeauthor{key}|        & only the authors \\
        \code|\citeyear{key}|          & only the year \\
        \bottomrule
    \end{tabular}
\end{table}

\code|\textcite| is for when the authors are the subject of the sentence,
\code|\parencite| for a citation at the end of a claim.
The square brackets are the optional argument from \autoref{sec:commands};
the tilde keeps \enquote{p.} and the number together.
With the style \texttt{numeric}, \code|\cite| and \code|\parencite| look the same;
with \texttt{authoryear}, \code|\cite| has no brackets.
\code|\textcite| and \code|\parencite| work as intended in every style, so prefer them.

\slidecode{04-citecmds.tex}
\sliderender{04-citecmds.pdf}

Four authors are shortened to \enquote{et al.} automatically,
and several sources share one pair of brackets, sorted: \enquote{[1, 2]}, not \enquote{[1] [2]}.

\code|\citeauthor| and \code|\citeyear| print a part of the citation, e.g.\ in \texttt{authoryear}:

\owncode{ch07-citemore.tex}
\ownrender{ch07-citemore.pdf}

\code|\footcite| puts the citation into a footnote,
which is how the citation styles of law and the humanities work (e.g.\ \texttt{style=verbose}).

\section{Citation styles}\label{sec:cite-styles}

The style is one word in the preamble.
The same document with \texttt{style=numeric}:

\slidecodetitled{04-biblatex-preamble.tex}{Preamble}
\slidecodetitled{04-styles.tex}{Document}
\sliderender{04-styles-numeric.pdf}

And with \texttt{style=authoryear}:

\slidecodetitled{04-styles-authoryear-preamble.tex}{Preamble}
\sliderender{04-styles-authoryear.pdf}

Every citation and the whole list change, and nothing in the document does.
With \texttt{authoryear}, the difference between \code|\textcite| and \code|\parencite| becomes visible:
the year moves into the brackets.

\texttt{style=alphabetic} uses short labels made from the names and the year:

\owncodetitled{ch07-styles-alphabetic-preamble.tex}{Preamble}
\ownrender{ch07-styles-alphabetic.pdf}

Many more styles exist as separate packages, e.g.\
\texttt{ieee} (package \texttt{biblatex-ieee}), common in engineering:

\owncodetitled{ch07-styles-ieee-preamble.tex}{Preamble}
\ownrender{ch07-styles-ieee.pdf}

\texttt{apa} (package \texttt{biblatex-apa}) is common in psychology and the social sciences,
and many journals and departments have their own.
Ask your supervisor which style to use;
it is one word to change.
